\newpage
\section{Introducción teórica}
Comenzaremos mencionando la importancia de la aplicación de los polinomios de Chebychev en el diseño de filtros.

\begin{itemize}
	\item \textbf{Características de atenuación específicas:} Los polinomios de Chebychev permiten diseñar filtros con características de atenuación específicas en la banda de paso. Esto significa que se pueden lograr atenuaciones más pronunciadas en la banda de paso en comparación con otros polinomios, como los de Butterworth.
	
	\item \textbf{Maximización de la banda de paso:} En el diseño de filtros, a menudo se desea maximizar la amplitud de la señal en la banda de paso mientras se garantiza una atenuación adecuada en la banda de rechazo. Los polinomios de Chebychev son ideales para este propósito, ya que minimizan la cantidad de atenuación en la banda de paso.
	
	\item \textbf{Permite diseñar filtros de orden bajo:} Los polinomios de Chebychev permiten diseñar filtros de orden bajo con características de atenuación específicas. Esto es beneficioso porque los filtros de orden bajo presentan más facilidad y menor costo para implementar.
	
	\item \textbf{Flexibilidad en el diseño:} La familia de polinomios de Chebychev es flexible y permite ajustar los parámetros del filtro, como la atenuación máxima en la banda de paso y la frecuencia de corte, para satisfacer los requisitos específicos del diseño.
	
	\item \textbf{Optimización de Ancho de Banda:} Los polinomios de Chebychev son útiles para optimizar el ancho de banda del filtro al mantener constante la tasa de atenuación en la banda de rechazo. Esto es esencial en aplicaciones donde el ancho de banda es crítico.
\end{itemize}